\chapter{Tests et validation}
\label{chap:validation}

\begin{chapeau}
Comment savoir si le logiciel donne le bon résultat ? Sans patient, nous avons testé sur des cas dont nous
connaissions la réponse : des déformités imposées à des os sains, un faux scanner de forme connue, et des campagnes de
centaines de cas contrôlés automatiquement. Ce chapitre présente la démarche et les résultats.
\end{chapeau}

\section{Trois niveaux de validation}

\fig{niveaux_validation.pdf}{Les trois niveaux de validation et leur état en novembre 2026.}{fig:niveaux}

Le niveau 1 vérifie que les calculs sont \textbf{justes} ; le niveau 2 qu'ils sont \textbf{robustes} ; le niveau 3, à
venir, qu'ils sont \textbf{utiles} au bloc opératoire. Les deux premiers niveaux ne disent rien de l'erreur finale
d'une vraie opération : elle viendra surtout des repères sur un os malade, de l'impression, de la pose du guide et du
geste du chirurgien.

\section{Les tests automatiques}

Le projet contient des tests qui se relancent en une commande (\cref{ann:commandes}).

\begin{table}[htbp]
\centering\small
\caption{Les tests automatiques du projet.}
\label{tab:tests}
\begin{tabularx}{\linewidth}{L{4.4cm}Y}
\toprule
\titrecol{Test} & \titrecol{Ce qu'il vérifie} \\
\midrule
\fichier{tests/test_chaine.py} & lecture et segmentation d'un faux scanner DICOM de forme connue : 14 contrôles (environ 3~s) \\
\fichier{tests/test_reperes.py} & repères sur l'os fictif posé au hasard et sur les os du lévrier : 40 contrôles (environ 13~s) \\
\fichier{tests/test_guide.py} & guides sur un os coudé synthétique (fermeture, ouverture avec dérotation, os voisin et cartilage) et sur les os du lévrier (environ 25~s) \\
\fichier{cora/tests/test_cas.mjs} & planification sur l'os fictif : 12 cas, 252 variantes et la torsion \\
\fichier{cora/tests/test_reel.mjs} & planification sur les os segmentés : tels quels, puis 6 déformités de test chacun \\
\fichier{tests/campagne_strategie.mjs} & 606 stratégies simulées, 9 règles \\
\fichier{tests/campagne_export.mjs} + \fichier{campagne_guides.py} & centaines de guides construits puis contrôlés « avec l'œil du chirurgien » \\
\bottomrule
\end{tabularx}
\end{table}

\subsection{Le faux scanner}

Le test \fichier{test_chaine.py} écrit en DICOM un fantôme dont on connaît tout : deux « os » en forme de capsule
(cylindre de 7~mm de rayon à bouts hémisphériques), bout à bout, séparés par un interligne de 1,5~mm et inclinés de
15\degre ; une corticale de 1,5~mm à 1300~HU, une moelle et des tissus mous à 40~HU, une table de scanner et un cylindre
de calibration ; un flou gaussien de 0,5~mm et un bruit de 40~HU. Les fichiers sont mélangés pour vérifier que le tri ne
dépend pas des noms. Le test vérifie la géométrie, la séparation des deux os, le rejet du fantôme de calibration,
l'étanchéité des maillages et la position de la surface (\cref{sec:demi-hauteur}).

\section{La segmentation}

\begin{itemize}
  \item \textbf{Position de la surface} (faux scanner) : biais de $+0{,}60$~mm avec le seul seuil, $+0{,}16$~mm après
  l'affinage à demi-hauteur ; objectif : moins de 0,5~mm.
  \item \textbf{Sensibilité au seuil} (lévrier, ± 50~HU) : 0,1 à 0,2~mm en moyenne, moins de 0,5~mm pour 95~\% des
  points (\cref{tab:sensibilite}).
  \item \textbf{Identification} : sur les trois scanners du lévrier, le fémur, le tibia, la fibula, la rotule, les
  fabelles, le talus, le calcanéus, le tarse, les métatarses, les phalanges et le bassin sont reconnus ; les maillages
  des os seuls sont étanches, de caractéristique d'Euler 2.
\end{itemize}

\section{Les repères}

Sur l'os fictif posé dans quatre positions au hasard (deux en miroir), le côté est reconnu 8 fois sur 8, les repères
sont à moins de 1~mm des repères exacts et identiques d'une pose à l'autre, et les angles de face diffèrent de moins de
1,3\degre{} de ceux du modèle. Sur le lévrier, les os coupés par le champ sont refusés et tous les angles des os
complets sont dans les normes (\cref{tab:normales-levrier}).

\section{La planification}

\begin{itemize}
  \item \textbf{Os fictif, 264 déformités} : l'angle de la déformité est retrouvé avec une erreur maximale de
  0,0065\degre.
  \item \textbf{Os du lévrier avec déformités de test} (référence : l'os avant la déformité) : erreur inférieure à
  0,001\degre ; après la correction simulée, les angles reviennent à ceux de l'os sain et une torsion de 15\degre{} est
  corrigée à 0,01\degre{} près ; les broches sont parallèles après la correction.
  \item \textbf{Référence « valeurs normales »} : l'écart naturel du lévrier à la moyenne publiée se retrouve dans la
  planification (plusieurs degrés). Ce n'est pas une erreur de calcul : quand c'est possible, il vaut mieux comparer au
  membre opposé du même chien.
\end{itemize}

\subsection{La campagne de stratégie}
\label{sec:campagne-strategie}

La campagne de stratégie (19/10/2026) planifie 606 cas : les 17 préréglages et leurs variantes (configurations de
broches, trait déplacé de ± 15~mm, charnière à 0,5 et 6~mm, chien en croissance) et 300 déformités tirées au hasard.
Pour chaque cas, les deux coins sont planifiés et contrôlés par 9 règles écrites à l'avance (\cref{tab:regles-r}).

\begin{table}[htbp]
\centering\small
\caption{Les 9 règles de la campagne de stratégie : 606 cas sur 606 sans erreur.}
\label{tab:regles-r}
\begin{tabularx}{\linewidth}{lYr}
\toprule
& \titrecol{Règle} & \titrecol{Erreurs} \\
\midrule
R1 & le guide n'est jamais du côté de la charnière & 0 \\
R2 & la charnière n'est dite conservée que si elle peut l'être (face du coin, règle 1, pas de dérotation, pont ≤ 3~mm) & 0 \\
R3 & la face d'appui est à 60\degre{} au plus de la face du coin, sinon l'abord non usuel est signalé & 0 \\
R4 & le coin recommandé est réalisable dès qu'un des deux l'est & 0 \\
R5 & le nombre et la disposition des broches sont ceux demandés & 0 \\
R6 & les broches sont parallèles après la correction & 0 \\
R7 & l'angle et la translation résiduels sont nuls & 0 \\
R8 & le niveau du trait et la profondeur de charnière demandés sont respectés & 0 \\
R9 & les angles après la correction sont ceux de l'os sain & 0 \\
\bottomrule
\end{tabularx}
\end{table}

Les coins recommandés suivent la pratique publiée : tibia en valgus → fermeture (84 cas) et tibia en varus → ouverture
(69 cas), avec l'accès comme critère décisif ; fémur → fermeture, avec la stabilité comme critère décisif.

\section{Les guides}

\subsection{La première campagne : un outil pour progresser}

La première campagne (14/10/2026) a construit et contrôlé des guides sur l'os fictif, avec ses os voisins et son
cartilage. Au début, seulement 46 cas sur 480 étaient sans erreur. Chaque type d'erreur a été corrigé dans le code, puis
la campagne a été relancée sur 640 cas, dont 320 tirés au hasard avec une graine jamais utilisée pendant le
développement : 640 sur 640 sans erreur (\cref{fig:campagne-guides}).

\fig{campagne_guides.png}{Les erreurs trouvées puis corrigées pendant la première campagne de guides.}{fig:campagne-guides}

Mesures après correction (minimum / médiane / maximum) : jeu guide--os 0,29 / 0,30 / 0,32~mm ; appui du guide 960 /
1470 / 2534~mm² ; guide à au moins 1~mm des os voisins et 1,5~mm du cartilage ; angles après la correction à
0,016\degre{} près ; broches parallèles à $2 \times 10^{-6}$ degré près.

\subsection{Les campagnes des guides v3 et du guide simplifié}

Après l'ajout d'une butée de profondeur, d'un verrou de réduction et de l'interdiction de l'espace entre les os, une
campagne (22/10/2026) a porté sur 520 cas : 320 cas « coins », où chaque paramètre est à ses valeurs extrêmes, et 200
cas tirés au hasard. Résultat : 520 sur 520 sans erreur. Le guide a ensuite été simplifié (06/11/2026, butée graduée
et verrou retirés, \cref{chap:guides}), et la même campagne a été relancée sur le nouveau guide : \textbf{520 sur 520
sans erreur} (\cref{tab:v3}).

\begin{table}[htbp]
\centering\small
\caption{La campagne du guide simplifié (06/11/2026). « Pas de guide » : coins que la planification déclare non
réalisables (elle conseille l'autre coin), construits quand même pour éprouver le programme ; le refus est le
comportement voulu.}
\label{tab:v3}
\begin{tabular}{lrr}
\toprule
& \titrecol{Coins (320)} & \titrecol{Hasard (200)} \\
\midrule
guide de coupe et guide de réduction construits & 266 & 159 \\
pas de guide : coin forcé, fibula sur le trajet de la lame & 54 & 41 \\
jeu guide--os mesuré & 0,296 à 0,327~mm & 0,298 à 0,324~mm \\
durée par guide (min / médiane / max) & 9 / 15 / 25~s & 10 / 15 / 23~s \\
\bottomrule
\end{tabular}
\end{table}

\subsection{Sur les os du lévrier}

Avec des déformités de test, les guides du fémur (coin de fermeture, torsion de test de 15\degre) et du tibia (coins de
fermeture et d'ouverture) sont construits sans alerte, avec un jeu de 0,31 à 0,33~mm, deux moitiés retirables et un
guide de réduction, en 10 à 16~s.

\section{La précision comparée}

\fig{precision.png}{L'erreur de notre calcul comparée à l'objectif du projet et aux résultats publiés.}{fig:precision}

L'erreur de calcul est si petite qu'elle ne se voit pas sur le graphique (\cref{fig:precision}). Les repères
automatiques (1,3\degre{} au plus de face sur l'os fictif) restent sous l'objectif de 2\degre. L'essentiel de l'erreur
d'une vraie opération viendra d'ailleurs.

\begin{aretenir}
\begin{itemize}
  \item Les calculs sont exacts à quelques millièmes de degré sur les cas simulés.
  \item Plus de 2 300 cas (stratégies et guides) passent tous les contrôles écrits à l'avance.
  \item Les campagnes ont surtout servi à \textbf{trouver des erreurs} : 434 cas sur 480 en avaient au départ.
  \item Rien n'a encore été imprimé ni essayé sur un os, et aucun vrai cas n'a été planifié : c'est le niveau 3.
\end{itemize}
\end{aretenir}
